\section*{Capítulo \ref{ch:b3:pericyclic} --- Reacciones pericíclicas}

\begin{solution}{exo:b3:pericyclic:1}
Diels--Alder: \omterm{def:b3:pericyclic:families}{cicloadición} [4+2]. Apertura del ciclobuteno: electrocíclica. Cope: \omterm{def:b3:pericyclic:families}{transposición sigmatrópica}
[3,3]. Pérdida de $\ce{SO2}$ del sulfoleno: queletrópica. Ozono y alqueno: \omterm{def:b3:pericyclic:dipolar}{cicloadición 1,3-dipolar}
([3+2]). Desplazamiento $[1,5]$-H: sigmatrópico.
\end{solution}

\begin{solution}{exo:b3:pericyclic:2}
Hexatrieno (6 electrones): térmico, \omterm{def:b3:pericyclic:rotation-modes}{disrotatorio}; fotoquímico, \omterm{def:b3:pericyclic:rotation-modes}{conrotatorio}.
Butadieno (4): térmico, \omterm{def:b3:pericyclic:rotation-modes}{conrotatorio}; fotoquímico, \omterm{def:b3:pericyclic:rotation-modes}{disrotatorio}.
\end{solution}

\begin{solution}{exo:b3:pericyclic:3}
Cuatro electrones, térmica: apertura \omterm{def:b3:pericyclic:rotation-modes}{conrotatoria}. Los dos grupos metilo, en caras
opuestas del anillo, giran ambos hacia fuera: (\textit{E},\textit{E})-hexa-2,4-dieno.
\end{solution}

\begin{solution}{exo:b3:pericyclic:4}
La reacción [2+2] \omterm{def:b3:pericyclic:faciality}{suprafacial} tiene $4q$ electrones: en el estado fundamental, la HOMO de un
eteno y la LUMO del otro se solapan con un extremo enlazante y otro antienlazante.
En el estado excitado, el $\pi^*$ simplemente ocupado de una de las moléculas desempeña el papel de la
HOMO, y ambos extremos coinciden.
\end{solution}

\begin{solution}{exo:b3:pericyclic:5}
Se conserva el \omterm{def:b3:point-groups:mirror}{plano de simetría}. Hexatrieno: $\psi_1$ S, $\psi_2$ A, $\psi_3$ S ocupados; $\psi_4$ A, $\psi_5$ S, $\psi_6$ A vacíos.
Ciclohexadieno, por orden: $\sigma$ S, $\pi_1$ S, $\pi_2$ A ocupados; $\pi_3^*$ S, $\pi_4^*$ A, $\sigma^*$ A vacíos. Uniéndolos por
simetría y por orden, $\psi_1 \to \sigma$, $\psi_2 \to \pi_2$, $\psi_3 \to \pi_1$: de ocupados a ocupados; el cierre \omterm{def:b3:pericyclic:rotation-modes}{disrotatorio}
está permitido por vía térmica.
\end{solution}

\begin{solution}{exo:b3:pericyclic:6}
Seis electrones con luz: \omterm{def:b3:pericyclic:rotation-modes}{conrotatorio}, de modo que los dos grupos metilo terminales acaban en
caras opuestas: \textit{trans}-5,6-dimetilciclohexa-1,3-dieno (la reacción térmica da
el isómero \textit{cis}).
\end{solution}

\begin{solution}{exo:b3:pericyclic:7}
El hidrógeno de C5 se desplaza a C1 del dieno, su vecino en el anillo,
a través de un estado de transición de seis electrones ($\sigma2_s + \pi4_s$), \omterm{def:b3:pericyclic:faciality}{suprafacialmente}: un primer desplazamiento
da el 1-metil-, y un segundo, el 2-metilciclopentadieno. Una cadena de cinco carbonos permite que el
hidrógeno se quede en una cara, cosa que el anillo pequeño impone de todos modos.
\end{solution}

\begin{solution}{exo:b3:pericyclic:8}
$\pi2_s$ (alqueno) $+ \pi2_a$ (C=C del ceteno, atacado por caras opuestas a través del
$\pi^*$ C=O perpendicular): el número de componentes $(4q + 2)_s$ es uno (el alqueno), impar:
permitida por vía térmica.
\end{solution}

\begin{solution}{exo:b3:pericyclic:9}
En la silla, el C=C del grupo vinilo (átomos 1--2), el oxígeno (3), el carbono del $\ce{OCH2}$
(4) y el C=C del alilo (5--6); el enlace O--C se rompe y se forma el enlace C1--C6:
el producto es el pent-4-enal,
$\ce{OHC-CH2-CH2-CH=CH2}$.
\end{solution}

\begin{solution}{exo:b3:pericyclic:10}
Ocho electrones son $4q$; con un componente \omterm{def:b3:pericyclic:faciality}{antarafacial}, la disposición es de topología de
Möbius, aromática con $4q$ electrones: el estado de transición es aromático y la
reacción está permitida por vía térmica (como el cierre \omterm{def:b3:pericyclic:rotation-modes}{conrotatorio} del octatetraeno).
\end{solution}

\begin{solution}{exo:b3:pericyclic:11}
Emparejar los mayores coeficientes une N3 con el carbono terminal: el
triazol 1,4-disustituido. Sin catalizador, los coeficientes difieren poco y
se forman los isómeros 1,4 y 1,5; el cobre(I) hace que la reacción transcurra por etapas a través de un
acetiluro de cobre y da únicamente el isómero 1,4.
\end{solution}

\begin{solution}{exo:b3:pericyclic:12}
El hidrógeno tiende un puente entre los extremos del SOMO de un radical de $j$ átomos, $k = (j + 1)/2$. Sus coeficientes
terminales son proporcionales a $\sin(k\pi/(j + 1))$ y $(-1)^{k+1}\sin(k\pi/(j + 1))$. Para $j = 7$, $k = 4$: signos opuestos, de modo que
los lóbulos del mismo signo están en caras opuestas en los dos extremos: \omterm{def:b3:pericyclic:faciality}{antarafacial}. Para
$j = 5$, $k = 3$: el mismo signo en la misma cara: \omterm{def:b3:pericyclic:faciality}{suprafacial}.
\end{solution}

\begin{solution}{pb:b3:pericyclic:1}
\textbf{1.} El enlace $\sigma$ C9--C10 del anillo B; con el 5,7-dieno da un hexatrieno,
C10=C5--C6=C7--C8=C9.

\textbf{2.} Seis electrones (los dos enlaces $\pi$ y el enlace $\sigma$): una apertura de anillo electrocíclica.

\textbf{3.} Por vía térmica, un ciclohexadieno y su hexatrieno abierto solo se equilibran
lentamente, a través de una barrera alta, y el anillo es el lado más estable; el
dieno absorbe la luz ultravioleta, y su estado excitado se abre en unos
picosegundos.

\textbf{4.} \omterm{def:b3:pericyclic:rotation-modes}{Conrotatoria}: en el estado excitado, la LUMO simplemente ocupada del sistema
trieno ($k = 4$) tiene extremos de signos opuestos.

\textbf{5.} El enlace C6=C7 procede del anillo: sus dos continuaciones de cadena, C5 y C8,
estaban del mismo lado de él, y siguen estándolo: \textit{Z}.

\textbf{6.} La apertura térmica, permitida solo de forma \omterm{def:b3:pericyclic:rotation-modes}{disrotatoria}, tiene una barrera muy superior a lo que
puede aportar el calor corporal, y conduciría cuesta arriba, hacia el trieno, menos estable.

\textbf{7.} De C19 (el metilo de C10) a C9: un desplazamiento sigmatrópico $[1,7]$-H.

\textbf{8.} Ocho: los seis electrones $\pi$ del trieno y los dos del enlace C--H.

\textbf{9.} $\sigma2 + \pi6$. \omterm{def:b3:pericyclic:faciality}{Suprafacialmente}, ambos serían s, con un componente $(4q + 2)_s$ cada uno:
par, prohibida. Con el trieno \omterm{def:b3:pericyclic:faciality}{antarafacial}, $\sigma2_s + \pi6_a$: un solo componente contado, impar,
permitida.

\textbf{10.} Siete átomos en una cadena helicoidal de configuración \textit{Z} permiten al hidrógeno pasar de una cara a
la otra; una cadena de tres átomos no puede torcerse tanto.

\textbf{11.} De Möbius (un componente \omterm{def:b3:pericyclic:faciality}{antarafacial}), con ocho electrones, $4q$: aromático.

\textbf{12.} Está permitido por vía térmica; su barrera se franquea lentamente a la temperatura corporal.

\textbf{13.} Seis electrones con luz: de nuevo \omterm{def:b3:pericyclic:rotation-modes}{conrotatorio}, pero puede girar los extremos en el
otro sentido y poner el H9 y el metilo C19 en las otras caras: un estereoisómero
de anillo cerrado, el lumisterol.

\textbf{14.} Con el enlace central \textit{E}, C19 y C9 quedan en lados opuestos de la cadena y muy
alejados: el hidrógeno no puede alcanzarlo.

\textbf{15.} La propia previtamina D$_3$ absorbe y se convierte en lumisterol y taquisterol,
que no dan vitamina D; se alcanza una mezcla fotoestacionaria, que limita la
producción.

\textbf{16.} La apertura ocurre en el estado excitado, en femtosegundos o picosegundos;
el desplazamiento es una reacción térmica con una barrera de \qty{85}{kJ/mol}.

\textbf{17.} $t_{1/2} = \ln 2/k = 0.693/\qty{0.023}{h^{-1}} = \qty{30}{h}$.

\textbf{18.} $1 - \eu^{-0.023 \times 24} = 0.42$.

\textbf{19.} $\ln 10/k = 2.303/\qty{0.023}{h^{-1}} = \qty{100}{h}$.

\textbf{20.} $k(\qty{293.15}{K}) = k(\qty{310.15}{K})\,\eu^{-(E_a/R)(1/293.15 - 1/310.15)} = 0.023 \times \eu^{-1.91} = \qty{3.4e-3}{h^{-1}}$.

\textbf{21.} $2.303/\qty{3.4e-3}{h^{-1}} = \qty{680}{h}$, unos 28 días.

\textbf{22.} La etapa térmica es siete veces más rápida a \qty{37}{\celsius} que a \qty{20}{\celsius}, de modo que la
previtamina formada en una mañana al sol se convierte en pocos días.

\textbf{23.} El esqueleto esteroideo, con sus numerosos estereocentros, viene ya hecho en
el esterol; la luz y el calor llevan a cabo después las dos etapas pericíclicas exactamente como en
la piel.

\textbf{24.} A la temperatura corporal, el 90\,\% de la previtamina D$_3$ se convierte en vitamina D$_3$ en unas
\qty{100}{h}, cuatro días.
\end{solution}
